\clearpage
\thispagestyle{empty}
\refstepcounter{section}
\label{app:electrical-worked-example}

\begin{center}
\textbf{APPENDIX C}

\vspace{\baselineskip}

\textbf{Worked Calculation for the Component 2\\
Electrical-Network Case}
\end{center}

\vspace{3\baselineskip}

\subsection{Purpose and Selected Scenario}

This appendix traces the reported Component 2 root-disturbance scenario with
\(B_{\mathrm{C2\_N01}}=0.75\), measured-loading scale \(1.00\), and neutral
transmission \(\tau=1\). It first demonstrates how phasor-derived apparent
power and declared capacity produce electrical susceptibility, then follows
the resulting risk through the 12-node observable arborescence. The inputs are
derived model-ready summaries; restricted raw phasor records and credentials
are not reproduced.

All nodes other than \texttt{C2\_N01} have \(B_i=0\) in this controlled
scenario. Because each non-root node in the observable Component 2 graph has
one predecessor, its update reduces to \(R_j=Q_{ij}=S_{ij}\tau_{ij}R_i\).
This simplification is specific to the selected single-source arborescence
case; it does not replace the general bounded aggregation rule.

\subsection{Step 1: Derive Baseline Electrical Susceptibility}

For the source transformer \texttt{tr\_117}, the 24 synchronized phasor
snapshots produced mean downstream real power \(P=985.9568\) kW and reactive
power \(Q=599.3408\) kvar. The loading is the mean of the apparent
powers computed at each snapshot, rather than the apparent power computed
from these two means. Writing the snapshot powers as \(P_t\) and \(Q_t\),

\[
\begin{aligned}
L_{12}
 &=\frac{1}{24}\sum_{t=1}^{24}\sqrt{P_t^2+Q_t^2}\\
 &=1153.8297\ \text{kVA}.
\end{aligned}
\eqtag{C-1}
\]

For comparison, \(\sqrt{\bar P^2+\bar Q^2}\approx1153.8285\) kVA.
The small difference arises from the order of averaging and the nonlinear
square-root operation; the snapshot-mean loading used below is unchanged.

\noindent\begin{minipage}{\textwidth}
Using the declared transformer capacity \(C_{12}=3750\) kVA, the source-edge
susceptibility is

\[
\begin{aligned}
S_{12}
 &=\min\left(1,\frac{L_{12}}{C_{12}}\right)\\
 &=\min\left(1,\frac{1153.8297}{3750}\right)\\
 &=0.3076879.
\end{aligned}
\eqtag{C-2}
\]
\end{minipage}

For comparison, \texttt{line\_408} has \(L=705.9977\) kVA and
\(C=1787.94\) kVA, giving

\[
S_{23}=\min\left(1,\frac{705.9977}{1787.94}\right)
=0.3948666.
\eqtag{C-3}
\]

The same rule is applied once to every observable edge. Table
\ref{tab:app-electrical-edge-inputs} lists the model-ready values needed for
the complete trace.

\begin{table}[H]
\caption{Component 2 Edge Inputs Used in the Worked Calculation}
\label{tab:app-electrical-edge-inputs}
\centering
\scriptsize
\begin{tabularx}{\textwidth}{@{}L L r r r@{}}
\toprule
Relation & Element & \(L\) (kVA) & \(C\) (kVA) & \(S=L/C\) \\
\midrule
C2\_N01 $\rightarrow$ C2\_N02 & \texttt{tr\_117} & 1153.8297 & 3750.00 & 0.3076879 \\
C2\_N02 $\rightarrow$ C2\_N03 & \texttt{line\_408} & 705.9977 & 1787.94 & 0.3948666 \\
C2\_N02 $\rightarrow$ C2\_N11 & \texttt{line\_409} & 375.4093 & 1787.94 & 0.2099675 \\
C2\_N03 $\rightarrow$ C2\_N04 & \texttt{line\_415} & 235.8730 & 1787.94 & 0.1319244 \\
C2\_N03 $\rightarrow$ C2\_N06 & \texttt{tr\_121} & 317.7301 & 1000.00 & 0.3177301 \\
C2\_N03 $\rightarrow$ C2\_N07 & \texttt{line\_417} & 153.8553 & 1787.94 & 0.0860517 \\
C2\_N03 $\rightarrow$ C2\_N09 & \texttt{line\_416} & 24.9112 & 1787.94 & 0.0139329 \\
C2\_N04 $\rightarrow$ C2\_N05 & \texttt{tr\_119} & 235.8730 & 1500.00 & 0.1572487 \\
C2\_N07 $\rightarrow$ C2\_N08 & \texttt{tr\_123} & 153.8553 & 750.00 & 0.2051404 \\
C2\_N09 $\rightarrow$ C2\_N10 & \texttt{tr\_125} & 24.9112 & 1000.00 & 0.0249112 \\
C2\_N11 $\rightarrow$ C2\_N12 & \texttt{tr\_127} & 375.4093 & 1500.00 & 0.2502728 \\
\bottomrule
\end{tabularx}
\end{table}

All listed ratios are below one, so clipping does not alter them in the
baseline scenario. The values are operating loading-to-capacity ratios, not
failure probabilities.

\subsection{Step 2: Initialize the Root and Evaluate the First Transfer}

The root has no predecessor, so the empty-product convention gives

\[
R_{\mathrm{C2\_N01}}=B_{\mathrm{C2\_N01}}=0.75.
\eqtag{C-4}
\]

The contribution through \texttt{tr\_117} is

\[
\begin{aligned}
Q_{12}
 &=S_{12}\tau_{12}R_1\\
 &=(0.3076879)(1.00)(0.75)\\
 &=0.2307659.
\end{aligned}
\eqtag{C-5}
\]

Because \(B_{\mathrm{C2\_N02}}=0\) and the node has only this predecessor,

\[
R_{\mathrm{C2\_N02}}
=1-(1-0)(1-0.2307659)=0.2307659.
\eqtag{C-6}
\]

\subsection{Step 3: Propagate Through the Main Branch}

Using the newly accumulated state at \texttt{C2\_N02}, the contribution to
\texttt{C2\_N03} is

\[
\begin{aligned}
Q_{23}
 &=(0.3948666)(1.00)(0.2307659)\\
 &=0.0911218,\\
R_{\mathrm{C2\_N03}}
 &=0.0911218.
\end{aligned}
\eqtag{C-7}
\]

The four outgoing branches from \texttt{C2\_N03} are then evaluated
independently. For example, the transformer branch to \texttt{C2\_N06}
gives

\[
\begin{aligned}
Q_{36}
 &=(0.3177301)(1.00)(0.0911218)\\
 &=0.0289521,\\
R_{\mathrm{C2\_N06}}
 &=0.0289521.
\end{aligned}
\eqtag{C-8}
\]

The parallel branch from \texttt{C2\_N02} to \texttt{C2\_N11} gives

\[
\begin{aligned}
R_{\mathrm{C2\_N11}}
 &=(0.2099675)(1.00)(0.2307659)\\
 &=0.0484533,\\
R_{\mathrm{C2\_N12}}
 &=(0.2502728)(1.00)(0.0484533)\\
 &=0.0121266.
\end{aligned}
\eqtag{C-9}
\]

These calculations show why risk decreases differently across branches: every
edge applies its own measured-loading susceptibility to the current upstream
state.

\subsection{Step 4: Complete the Topological Trace}

Table~\ref{tab:app-electrical-trace} reports the complete one-pass trace. Minor
differences in the displayed products can result from rounding; the evaluator
uses the unrounded input values listed in the model tables.

\begin{table}[H]
\caption{Component 2 Root-Severity Trace at \(B_{N01}=0.75\)}
\label{tab:app-electrical-trace}
\centering
\small
\begin{tabularx}{\textwidth}{@{}c L L r@{}}
\toprule
Depth & Node & Calculation using unrounded predecessor state & Final \(R_i\) \\
\midrule
0 & C2\_N01 & Local disturbance & 0.7500000 \\
1 & C2\_N02 & \(0.75(0.3076879)\) & 0.2307659 \\
2 & C2\_N03 & \(R_{N02}(0.3948666)\) & 0.0911218 \\
2 & C2\_N11 & \(R_{N02}(0.2099675)\) & 0.0484533 \\
3 & C2\_N04 & \(R_{N03}(0.1319244)\) & 0.0120212 \\
3 & C2\_N06 & \(R_{N03}(0.3177301)\) & 0.0289521 \\
3 & C2\_N07 & \(R_{N03}(0.0860517)\) & 0.0078412 \\
3 & C2\_N09 & \(R_{N03}(0.0139329)\) & 0.0012696 \\
3 & C2\_N12 & \(R_{N11}(0.2502728)\) & 0.0121266 \\
4 & C2\_N05 & \(R_{N04}(0.1572487)\) & 0.0018903 \\
4 & C2\_N08 & \(R_{N07}(0.2051404)\) & 0.0016085 \\
4 & C2\_N10 & \(R_{N09}(0.0249112)\) & 0.0000316 \\
\bottomrule
\end{tabularx}
\end{table}

Summing the unrounded states gives

\[
\sum_{i=1}^{12}R_i=1.1860821582\approx1.186,
\eqtag{C-10}
\]

which reproduces the risk sum reported in Chapter IV. The three nodes above
the declared 0.05 affected-node threshold are \texttt{C2\_N01},
\texttt{C2\_N02}, and \texttt{C2\_N03}; the maximum terminal-node state is
\(R_{\mathrm{C2\_N06}}=0.0289521\approx0.029\).

The calculation establishes that measured normal-operation loading can be
converted into heterogeneous bounded susceptibilities and evaluated by the
unchanged GBBRPM core. The injected \(B=0.75\) is controlled rather than an
observed electrical fault, so the result is mechanism and sensitivity
evidence rather than event-outcome validation.
